\documentclass[11pt]{article}
\usepackage[margin=1in]{geometry}
\usepackage{amsmath,amssymb,amsthm}
\usepackage[T1]{fontenc}
\usepackage{lmodern}
\usepackage[hidelinks]{hyperref}
\newtheorem{theorem}{Theorem}
\title{The missing factor in Conjecture 00000003454}
\author{gaochengzhecpu}
\date{October 9, 2026}
\begin{document}
\maketitle
\begin{abstract}
For the smooth polynomial $f(x)=(x-1/2)^2$, the Bernstein approximation
at $x=1/2$ has error exactly $1/(4n)$. Thus the error multiplied by $n$
converges to $1/4$, whereas $f''(1/2)=2$. This disproves the source's
identification of the deviation limit with the second derivative.
The unscaled error also has a different limit, namely zero.
\end{abstract}

\section{Statement and interpretation}
The source conjecture describes the Voronovskaya asymptotics of Bernstein
operators by saying that ``the limit of the deviation is the second
derivative,'' with a rate inverse in $n$. The standard Bernstein operator is
\[
 (B_nf)(x)=\sum_{k=0}^{n}f(k/n)\binom nk x^k(1-x)^{n-k},\qquad n\geq1.
\]
We examine the usual normalized deviation $n(B_nf-f)$ and, separately,
the literal unscaled deviation $B_nf-f$. Neither converges to $f''$ in
our example, even at an interior point and for a polynomial.
The issue is the coefficient of the second derivative, not the occurrence
of an order $1/n$ approximation error.

\section{Exact counterexample}
\begin{theorem}
For $f(x)=(x-1/2)^2$ and every positive integer $n$,
\[
 (B_nf)(1/2)-f(1/2)=\frac{1}{4n}.
\]
Consequently, $n((B_nf)(1/2)-f(1/2))$ converges to $1/4$, and the
unscaled deviation converges to $0$. Both limits differ from $f''(1/2)=2$.
\end{theorem}
\begin{proof}
The binomial variance identity gives, for $0\leq x\leq1$,
\[
 \sum_{k=0}^{n}(k/n-x)^2\binom nk x^k(1-x)^{n-k}
   =\frac{x(1-x)}n.
\]
At $x=1/2$, the left side is exactly $(B_nf)(1/2)$, while the right side
is $1/(4n)$. Since $f(1/2)=0$, this proves the displayed formula.
Multiplying by $n$ gives the constant $1/4$. Without that multiplication,
the error tends to zero. Direct differentiation gives
$f'(x)=2(x-1/2)$ and $f''(x)=2$. The asserted limit is therefore false
under either reading of the deviation.
\end{proof}

\paragraph{The coefficient matters.}
The usual Voronovskaya expression contains the factor $x(1-x)/2$ in front
of $f''(x)$; at the midpoint this factor is $1/8$ and gives $1/4$ for our
quadratic. The proof above is an exact finite-sum calculation and does not
assume a general Voronovskaya theorem. It does not dispute the classical
formula with its factor restored, or any separately specified
renormalization that divides out this factor. Neither such factor nor
such a renormalization appears in the source.

\section{Formal verification}
The Lean project uses Mathlib's actual \texttt{bernsteinApproximation}
on continuous functions on $[0,1]$, its proved binomial variance identity,
real derivatives, and filter limits. In particular:
\begingroup\raggedright
\begin{itemize}
\item \texttt{quadratic\_\allowbreak smooth} proves the polynomial is infinitely
      differentiable; \texttt{quadratic\_\allowbreak second\_\allowbreak derivative} evaluates
      the second derivative at every real point.
\item \texttt{approximation\_\allowbreak at\_\allowbreak midpoint} identifies the actual
      Bernstein sum with $1/(4n)$ for every positive integer $n$.
\item \texttt{scaled\_\allowbreak deviation\_\allowbreak limit} and \texttt{deviation\_\allowbreak limit}
      prove the normalized and unscaled limits. The sequences are indexed
      by $n+1$ so every approximation order is positive.
\item \texttt{scaled\_\allowbreak second\_\allowbreak derivative\_\allowbreak limit\_\allowbreak false} and
      \texttt{raw\_\allowbreak second\_\allowbreak derivative\_\allowbreak limit\_\allowbreak false} state the two
      contradictions using the actual second derivative, not a
      substituted numerical surrogate.
\end{itemize}
\endgroup
No auxiliary numerical computation is required. All final theorems use
only Lean's standard logical axioms, with no unproved placeholders,
custom axioms, or native evaluation used as a proof.

\paragraph{Reproduction.}
In \texttt{lean/}, run:
\begin{verbatim}
lake build
lake env lean -DwarningAsError=true Main.lean
\end{verbatim}
The project pins Lean 4.19.0 and an exact Mathlib revision.
Run \texttt{tectonic main.tex} to rebuild the PDF.
\texttt{SOURCE.md} preserves the bilingual conjecture.

\begingroup\raggedright
\begin{thebibliography}{9}
\bibitem{source} The Last Math Competition,
\href{https://github.com/The-Last-Math-Competition/The-Last-Math-Competition/blob/main/conjectures/00000003454.md}{Conjecture 00000003454}.
\bibitem{mathlib} The Mathlib Community, Lean 4 mathematical library,
\href{https://github.com/leanprover-community/mathlib4/commit/c44e0c8ee63ca166450922a373c7409c5d26b00b}{revision c44e0c8ee63c},
module \texttt{Analysis.SpecialFunctions.Bernstein}.
\end{thebibliography}
\endgroup
\end{document}
